\documentclass[11pt]{article}

% ---- arXiv 스타일 프리프린트 (XeLaTeX + 한글) ----
\usepackage[margin=1in]{geometry}
\usepackage{amsmath,amssymb}
\usepackage{graphicx}
\usepackage{booktabs}
\usepackage{array}
\usepackage{multirow}
\usepackage{longtable}
\usepackage{caption}
\usepackage{authblk}
\usepackage{xcolor}
\usepackage{enumitem}
\usepackage{textcomp}
\usepackage{mhchem}
\usepackage{siunitx}
\sisetup{detect-all}
\usepackage{pgfplots}
\pgfplotsset{compat=1.15}
\usetikzlibrary{arrows.meta,positioning,calc}
\usepackage{kotex}
\usepackage{fontspec}
\setmainhangulfont{NanumMyeongjo}
\setsanshangulfont{NanumGothic}
\usepackage[hidelinks]{hyperref}

\captionsetup{font=small,labelfont=bf}
\renewcommand{\arraystretch}{1.15}

% 편의 매크로
\newcommand{\dpH}{\ensuremath{\Delta\mathrm{pH}}}
\newcommand{\dkh}{\ensuremath{\,\mathrm{dKH}}}
\newcommand{\co}{\ensuremath{\mathrm{CO_2}}}
\newcommand{\pco}{\ensuremath{p_{\mathrm{CO_2}}}}
\newcommand{\at}{\ensuremath{A_\mathrm{T}}}

\title{\bfseries 해수 탄산알칼리니티 측정을 위한 저가형 단일프로브\\
차동 pH 폭기 모니터: 계통오차 분해와 밀폐 공통 헤드스페이스\\
평형화를 통한 0.25\,dKH 이하 정확도 달성}

\author[1]{nowforyou\thanks{교신: \texttt{taeseokyi@gmail.com}}}
\affil[1]{REEF FAMILY}

\date{\today}

\begin{document}
\maketitle

\begin{abstract}
\noindent
탄산알칼리니티(이하 ``KH'', 단위 \dkh)는 산호 어항에서 가장 제어가 중요한
수질 인자이지만, 저렴한 연속 모니터는 드물고 시약 적정은 번거롭다. 본 논문은
시료를 대기 \co{}와 평형에 이를 때까지 폭기한 뒤 도달한 평형 pH로부터
알칼리니티를 추정하는 무시약(reagent-free) 개방형 모니터를 기술한다. \emph{단일}
pH 프로브로 어항 시료(``tank'')와 알칼리니티가 알려진 참조수(``ref'')를 연달아
측정한다. dKH가 두 pH의 \emph{차}로 계산되므로
($\mathrm{KH_{tank}} = \mathrm{KH_{ref}}\cdot 10^{\dpH}$), 프로브의 절대
드리프트와 실내 \co{} 변동은 공통모드 신호로서 상쇄된다. 9일간의 계측 캠페인
(5개 펌웨어/방법 버전에 걸친 29회 측정)에서 무보정 오차 $-0.92\dkh$가 지속적으로
관찰되었으며, 우리는 이를 영점 보정으로 지우는 대신 \emph{분해}하였다. 세 가지
개입---(i) 진짜 가스 평형까지 가는 평탄 추종 측정, (ii) 시료ㆍ참조수ㆍ펌프 흡입이
공유하는 단일 \emph{밀폐 공통 헤드스페이스}, (iii) 헤드스페이스 \pco{}를 고정하는
능동 폭기 5\,L 참조수 앵커---이 새벽 \co{} 오차와 일주기 증발 교란을 붕괴시켰고
측정 산포를 $\mathrm{SD}=\pm0.08\dkh$ (목표 $\pm0.25\dkh$의 약 1/3)로 줄였다.
이어진 널 실험(두 용기에 동일한 물, 참값 $\dpH=0$)이 잔차를 국소화하였다:
\emph{완전 널}에서 $\dpH=-0.003$, $\mathrm{dKH}=8.387$, 즉 계측 오프셋은 단
$-0.06\dkh$에 불과했다. 남은 $-0.86\dkh$는 수일 된, 소부피, 개방 상태의 tank
분취액에서 일어난 알칼리니티 손실(\co{} 과포화에 의한 \ce{CaCO3} 석출/생물막)로
추적되었다---이는 측정 결함이 아니라 시료 취급 아티팩트다. 같은 노화 분취액으로 하루
뒤 널을 재측정하니 전자장치를 건드리지 않았는데도 널이 $-0.06$에서 $-0.22\dkh$로
무너지는 것을 관찰하여, 이 귀속을 시간 영역에서 확증하였다. 따라서 본 장치는 영점
보정 없이 $\pm0.25\dkh$ 목표를 달성하며, 영점 보정은 불필요한 것으로 판정되었다.
우리가 독립적으로 도출한 오차 구조(실내 \co{} 민감도, 기포 품질 의존성, 해법으로서의
밀폐)는 상용 제품의 보고된 거동을 그대로 재현하여, 진단과 해법 양쪽에 외부 검증을
제공한다.
\end{abstract}

\noindent\textbf{키워드:} 산호 어항; 탄산알칼리니티; 총알칼리니티; \co{} 평형화;
차동 pH; 저가 계측; 공통모드 제거; \ce{CaCO3} 석출.

% =====================================================================
\section{서론}
% =====================================================================
경산호를 비롯한 석회화 생물은 탄산칼슘을 침착시키며 탄산알칼리니티를 끊임없이
소모한다. 닫힌 산호 어항에서 이 알칼리니티 감소는 빠르며(흔히
$0.5$--$1.5\dkh\,\mathrm{day^{-1}}$), $\sim1\dkh$를 넘는 변동은 산호 조직 스트레스와
연관된다. 알칼리니티(\at{}, 취미가들은 보통 독일식 도 단위의 ``KH''로 표기,
$1\dkh \approx 0.357\,\mathrm{meq\,L^{-1}}$)는 따라서 \emph{자주} 측정할 가치가
가장 큰 인자다.

접근 가능한 선택지는 연속 사용에 부족하다. 수동 시약 적정(예: 상용 드롭 키트)은
정확하지만 노동집약적이고 이산적이다. 색도 로봇은 적정을 자동화하지만 시약을
소모하고 비싸다. 직접 전위차 pH 단독은 \emph{알칼리니티 측정이 아니다}. 해수 pH는
\at{}와 \pco{}에 함께 의존하고, 어항 \pco{}는 하루 동안 크게 변하기 때문이다(실내
환기, 광합성, 호흡).

무시약 대안은 다음 사실을 이용한다: 시료를 \pco{}가 주변 공기와 평형에 이를 때까지
폭기하면, 평형 pH는 고정된 온도ㆍ염분에서 알칼리니티의 일가 함수가 된다. 상용
제품(AquaWiz)이 이 원리로 만들어져 취미가들이 사용하지만, 커뮤니티 보고는 실내
\co{} 변동과 폭기/기포 품질에 결부된 반복적 부정확성을 기록한다. 본 논문의 기여는
원리 자체가 아니라 \emph{정량적 진단}이다: 작동하는 DIY 제작물에서 출발해 9일
캠페인을 계측하고, $-0.92\dkh$의 원시 오차를 드러내며, 널 실험으로써(구성상) 그
거의 전부가 장치 한계가 아니라 시료 취급 아티팩트임을 보인다---진짜 계측 오프셋을
$-0.06\dkh$로 줄이고 무보정으로 $\pm0.25\dkh$ 목표를 달성한다. 그 과정에서 임의의
단일프로브 차동 pH 알칼리니티 모니터에 일반화될 것으로 기대되는 일련의 설계 규칙
(밀폐 공통 헤드스페이스; 진짜 평형까지의 평탄 추종; 능동 폭기 참조수 앵커; 이중
게이트 평탄 판정)을 확립한다.

% =====================================================================
\section{측정 원리}\label{sec:principle}
% =====================================================================

\subsection{폭기가 평형 pH를 알칼리니티의 프록시로 만든다}
고정된 온도 $T$와 염분 $S$에서 해수 탄산계는 자유도가 둘이다;
$\{\at{}, \mathrm{DIC}, \pco{}, \mathrm{pH}\}$ 중 둘을 고정하면 나머지가 정해진다.
실내 공기를 과량으로 폭기하면 용존 \co{}가 기상과의 헨리 법칙 평형
$[\co{}]_{aq}=K_H\,\pco{}$로 구동되어 $\pco{}(\mathrm{water}) \to
\pco{}(\mathrm{air})$가 된다. \pco{}가 대기에 의해 고정되면, 평형 pH는
알칼리니티의 일가 함수가 된다:
\begin{equation}
\mathrm{pH_{eq}} = f(\at{},\,\pco{},\,T,\,S).
\label{eq:pheq}
\end{equation}
\pco{}, $T$, $S$가 공통이면 $\mathrm{pH_{eq}}$를 움직일 수 있는 \emph{유일한}
변수는 \at{}이다. 따라서 측정된 평형 pH는 알칼리니티의 프록시가 된다.

\subsection{차동이 실내 \co{}를 상쇄하는 이유}
본 장치는 \pco{}의 절대값에 결코 의존하지 않는다. 참조수와 tank 시료를 \emph{같은}
대기 천장까지 폭기하고(최종 설계에서는 하나의 공유 헤드스페이스에 대해 둘을 동시
폭기, \ref{sec:methods}절), 두 평형 pH의 \emph{차}를 보고한다. 절대 \pco{} 항이
두 시료에 공통이므로 차에서 상쇄된다: 실내 \co{} 변동은 떨어져 나가는 공통모드
신호다. 차감 후 남는 것은 진짜 알칼리니티 차에 전극 잡음($\pm0.005$--$0.01$\,pH)을
더한 것이다. 이것이 이 방법의 핵심 견고성 논거이며, (두 개의 정합 프로브가 아닌)
\emph{단일} 프로브를 쓰는 이유다: 단일 센서의 느린 비대칭 전위와 기준액 접합부
드리프트는 두 측정에 동일하게 실려 \dpH{}에서 상쇄된다.

\subsection{과폭기는 불가능하고, 오직 \emph{불충분} 폭기만이 오차를 낳는다}
식~\eqref{eq:pheq}는 $\pco{}(\mathrm{air})$에서 단단한 천장을 갖는다: 물이 대기
평형에 도달하면 추가 폭기는 pH를 더 올릴 수 없다. 따라서 ``과폭기'' 실패 모드는
존재하지 않는다. 오차는 오직 \emph{불충분}을 통해 들어온다: 한 용기에 잔류 \co{}가
남으면 그쪽 pH가 부당하게 눌려 거짓 \dpH{}가 나타난다. 차동의 전제는 바로 ref와
tank가 같은 \pco{}를 공유한다는 것이다; 대칭이 깨지면 상쇄가 깨진다. 이것이 고정
시간이 아니라 진짜 평탄까지 측정하는 동기다(\ref{sec:methods}절).

\subsection{완전 평형에서의 부피 무관성}
식~\eqref{eq:pheq}에는 부피 항이 없다(헨리 법칙은 시강성질). 따라서 큰 부피
불일치---최종 널에서 $100\,\mathrm{mL}$ tank 물 대 $5\,\mathrm{L}$ 참조수---조차
두 용기가 \emph{완전} 평형에 도달하기만 하면 차동 오차는 \emph{0}이다. 익숙한
``대칭을 위한 등부피'' 규칙은 부피가 다르면 일차 접근 곡선
$C(t)-C^\ast = (C_0-C^\ast)\,e^{-t/\tau}$의 다른 지점에 놓이는 \emph{부분} 평형에서만
적용된다. 진짜 요건은 단지 (a) 둘 다 완전 평형, (b) \pco{}와 $T$ 정합뿐이다. 더 큰
부피가 율속이며, 시간상수는
\begin{equation}
\tau \;\propto\; \frac{V}{K_H\,Q_\mathrm{gas}}\,\beta
\;\sim\; \frac{1}{\mathrm{VVM}}\,\beta,
\label{eq:tau}
\end{equation}
이다. 여기서 $\mathrm{VVM}=Q_\mathrm{gas}/V$는 단위 액 부피당 분당 기체 부피이고
$\beta$는 해수 화학 인자다(미량 $[\co{}]_{aq}$만 탈기 가능하며 Revelle 완충이 $\tau$를
$\sim8$--$15\times$ 키운다). 식~\eqref{eq:tau}는 뒤에서 논의할 ``폭기를 키워 측정을
단축'' 레버와, 율속이 되는 참조수를 소형화한 선택의 근거다.

\subsection{차동 dKH 산출식}
펌웨어(Arduino Nano / ATmega328P; pH는 DFRobot Gravity 프로브 + ADS1115 16비트 ADC)는
단일 프로브를 디지털화하여 참조수, 이어서 tank pH를 읽고 다음을 계산한다:
\begin{equation}
\boxed{\;\mathrm{KH_{tank}} \;=\; \mathrm{KH_{ref}}\cdot 10^{-(\mathrm{pH_{ref}}-\mathrm{pH_{tank}})}
\;=\; \mathrm{KH_{ref}}\cdot 10^{\dpH}\;}
\label{eq:dkh}
\end{equation}
여기서 $\dpH \equiv \mathrm{pH_{tank}}-\mathrm{pH_{ref}}$이고
$\mathrm{KH_{ref}}=8.448\dkh$는 EEPROM에 저장된 \emph{앵커} 상수다(증발 보충 후
\texttt{calref} 명령으로 설정; 매 측정값이 \emph{아님}---이 때문에 로그의 참조수
열은 항상 상수 8.448로 찍힌다). 식~\eqref{eq:dkh}는 \dpH{}에만 의존하며 절대 pH는
정확히 상쇄된다. 미분하면,
\begin{equation}
\frac{d\,\mathrm{KH_{tank}}}{d\,\dpH}\bigg|_{\dpH=0}
= \mathrm{KH_{ref}}\ln 10 \approx 19.45\;\dkh\,\mathrm{pH^{-1}},
\label{eq:sens}
\end{equation}
즉 상쇄되지 못한 $0.01$\,pH 차가 $\approx0.19\dkh$로 사상된다. 이 높은 민감도가
상쇄와 평탄 판정 둘 다를 중요하게 만든다. 온도 보정은 매 측정마다 Nernst 관계
$\mathrm{pH_{corr}}=7.0+(\mathrm{pH_{raw}}-7.0)\,(273.15+T_\mathrm{cal})/(273.15+T_\mathrm{meas})$로
적용되며, 15--35\,\textcelsius에서 명목상 0 오차다. 각 pH 읽기는 64 샘플의
절사평균(상하 16개씩 버리고 가운데 32개 평균)으로 폭기 기포 스파이크를 배제한다.

% =====================================================================
\section{장치}\label{sec:apparatus}
% =====================================================================
제작물(표~\ref{tab:bom}, 그림~\ref{fig:device})은 4개의 양방향 연동펌프, 2개의 USB
에어펌프, 단일 pH 프로브, 1-Wire 온도 센서를 둘러싼 폐쇄 유체 루프이며, Arduino
Nano가 9600\,baud 블루투스(HC-06) 시리얼로 호스트 PC와 통신한다.

\begin{figure}[t]
\centering
\begin{tikzpicture}[
  font=\footnotesize, >={Latex[length=1.8mm]},
  vessel/.style={draw, rounded corners=2pt, fill=blue!6, align=center},
  pump/.style={draw, circle, fill=yellow!35, font=\scriptsize\bfseries, inner sep=0.4pt, minimum size=5.5mm},
  flow/.style={<->, thick, blue!55!black},
  air/.style={-{Latex[length=1.6mm]}, densely dotted, red!65!black, thick},
  wire/.style={-{Latex[length=1.6mm]}, gray!55!black},
]
\draw[dashed, rounded corners=7pt, gray!65, line width=0.9pt] (0.1,-3.05) rectangle (10.9,1.95);
\node[anchor=north east, gray!55!black, font=\scriptsize] at (10.8,1.88) {밀폐 공통 헤드스페이스 ($\approx$4\,L)};
\node[draw, rounded corners=2pt, fill=gray!10, align=center] (mcu) at (4.9,2.75)
  {Arduino Nano $+$ HC-06\,/\,ADS1115 pH\,/\,DS18B20};
\node[vessel, minimum width=15mm, minimum height=20mm, fill=cyan!10] (aq) at (-2.2,-0.6) {본수조\\(어항)};
\node[vessel, minimum width=14mm, minimum height=14mm] (hold) at (1.9,0.7) {홀딩};
\node[vessel, minimum width=21mm, minimum height=21mm] (meas) at (4.9,-0.7) {측정\\챔버};
\node[vessel, minimum width=16mm, minimum height=19mm] (ref) at (8.9,-0.4) {참조수\\1\,L\\\scriptsize($+$온도)};
\node[vessel, minimum width=15mm, minimum height=11mm, fill=green!8] (kcl) at (4.9,-2.45) {3\,M KCl};
\node[draw, rounded corners=1pt, fill=red!9, font=\scriptsize, align=center] (airp) at (8.9,-2.45) {에어펌프 (\texttt{ron})};
\draw[line width=1pt] ([xshift=-3mm]meas.north) -- ++(0,5mm) coordinate (pp);
\node[font=\scriptsize, anchor=south] at (pp) {pH 프로브};
\draw[flow] (aq.east) -- (hold.west);
\node[pump] at ($(aq.east)!0.5!(hold.west)$) {M1};
\draw[flow] (hold.south) -- (meas.north west);
\node[pump] at ($(hold.south)!0.5!(meas.north west)$) {M2};
\draw[flow] (ref.west) -- (meas.east);
\node[pump] at ($(ref.west)!0.5!(meas.east)$) {M4};
\draw[flow] (kcl.north) -- (meas.south);
\node[pump] at ($(kcl.north)!0.5!(meas.south)$) {M3};
\draw[air] (airp.north) -- (ref.south);
\draw[air] (airp.west) to[out=180,in=-90] (meas.south east);
\draw[wire] (pp) to[out=90,in=-90] (mcu.south);
\end{tikzpicture}
\caption{밀폐 폐루프 모니터 개념도. 측정 챔버의 단일 pH 프로브가 tank 분취액과
참조수를 차례로 읽는다; 4개의 양방향 연동펌프(M1--M4)가 본수조ㆍ홀딩ㆍ측정 챔버ㆍ
참조수 저장조ㆍKCl 보관 사이로 액체를 옮긴다; 2개의 에어펌프(\texttt{ron})가 참조수와
챔버를 동시 폭기한다. 참조수ㆍ챔버ㆍ펌프 흡입이 하나의 밀폐 헤드스페이스($\approx$4\,L)를
공유하여 시료와 참조수가 동일한 $p_{\mathrm{CO_2}}$를 보고 증발이 억제된다; 컨트롤러는
바깥에 있다. 측정 사이 프로브는 3\,M KCl(M3)에 담겨 보관된다.}
\label{fig:device}
\end{figure}

\begin{table}[t]
\centering
\caption{부품 목록(현재 제작물).}
\label{tab:bom}
\footnotesize
\begin{tabular}{@{}p{0.26\linewidth} p{0.50\linewidth} c@{}}
\toprule
서브시스템 & 부품 & 수량 \\
\midrule
마이크로컨트롤러 & Arduino Nano V3.0 (ATmega328P) & 1 \\
pH 프로브 & DFRobot Gravity pH V2 (SEN0161-V2), 단일 유리전극 & 1 \\
ADC & DFRobot Gravity I\textsuperscript{2}C ADS1115 (16비트) & 1 \\
온도 & DS18B20 방수(1-Wire), 5\,L 참조수에 삽입 & 1 \\
통신 & HC-06 블루투스, 9600\,baud & 1 \\
펌프 & NKP-DC-B06B 12\,V 연동, 양방향 (M1--M4) & 4 \\
모터 드라이버 & L298 2\,A 모듈(펌프 2 + 에어 1) & 3 \\
에어펌프 & USB 5\,V, 동시 스위칭(\texttt{ron}) & 2 \\
참조수 저장조 & 5\,L 밀폐 용기(``위즈 수조'') & 1 \\
챔버 & 200\,mL 비커: 홀딩, 측정, 3\,M KCl & 3 \\
프로브 보관 & 3\,M KCl ($\approx$22\,\% w/v) & --- \\
\bottomrule
\end{tabular}
\end{table}

\paragraph{유체 배관 토폴로지(현재).} 모든 펌프는 양방향이다($\mathtt{f}$=정방향,
$\mathtt{b}$=역방향):
\begin{itemize}[nosep,leftmargin=1.4em]
\item \textbf{M1}: 본수조 $\leftrightarrow$ \emph{홀딩} 비커;
\item \textbf{M2}: 홀딩 $\leftrightarrow$ \emph{측정} 챔버;
\item \textbf{M3}: 3\,M KCl 비커 $\leftrightarrow$ 측정 챔버(프로브 소크/보관);
\item \textbf{M4}: 5\,L 참조수 $\leftrightarrow$ 측정 챔버.
\end{itemize}
tank 시료를 중간 홀딩 비커로 경유시키면 참조수 측정 동안 tank를 임시
\emph{파킹}할 수 있어 위상 간 지연을 최소화한다(\ref{sec:methods}절). 가동 시간은
$\mathtt{mXf{:}NN}$ / $\mathtt{mXb{:}NN}$의 정수(초)로 인코딩되며, 호스를 완전히
비우기 위해 역방향이 정방향보다 $\sim$8--10\,s 길다(현재 값
$\mathtt{m1f{:}70/m1b{:}82}$, $\mathtt{m2f{:}60/m2b{:}68}$,
$\mathtt{m3f{:}60/m3b{:}68}$, $\mathtt{m4f{:}60/m4b{:}70}$; 5\,L 호스가 가장 길어
\texttt{m4b:70}).

\paragraph{제어 명령.} 에어와 모터 인에이블은 독립 라인이다: \texttt{ron}은 \emph{두}
에어펌프를 켠다(액체가 정지한 동안에만 사용---기포 청소와 측정); \texttt{ton}은 PWM
모터 제어기를 켠다(어떤 펌프든 돌기 전에 \emph{필수}); \texttt{airoff}는 둘 다 끈다.
\texttt{airoff}가 PWM 제어기도 끄므로, 모든 액체 이동은
\texttt{airoff}\,$\rightarrow$\,\texttt{ton} 순으로 선행되어야 한다는 핵심 불변식이
따른다; 모터 직전의 마지막 에어 명령은 반드시 \texttt{ton}이어야 한다. 측정
프리미티브는 \texttt{tank} / \texttt{ref}(한 pH 읽기 샘플링ㆍ저장)와
\texttt{calkh}(식~\eqref{eq:dkh} 적용)이다.

\paragraph{밀폐 공통 헤드스페이스(필수).} 5\,L 참조수, 측정 챔버, 펌프 흡입이
\emph{하나의} 밀폐 헤드스페이스를 공유한다. 이는 선택적 개선이 아니다: (i) 측정
순간 tank와 ref가 동일한 \pco{}를 보게 보장하고(차동의 전제), (ii) 실내 \co{} 변동을
원천 차단하며, (iii) 증발을 억제한다. 최종 구성에서 헤드스페이스는 $\approx$4\,L,
참조수 $\approx$1\,L, tank 분취액 $100\,\mathrm{mL}$이며; 참조수 \emph{안에} 둔
에어스톤이 공유 \pco{}를 수 분 내 능동적으로 고정한다(``앵커'',
\ref{sec:methods}절). 분리되거나 격리된 헤드스페이스는 명시적으로 금지된다---두
헤드스페이스는 독립적으로 드리프트하여 상쇄를 파괴한다.

% =====================================================================
\section{방법}\label{sec:methods}
% =====================================================================

\subsection{2-위상 평탄 추종 사이클(V4)}
평탄 판정은 호스트 측 결정 루프이므로 PC 스크립트가 펌웨어를 단계별로 구동한다.
사이클은 \textbf{tank 먼저, 그다음 ref}를 측정하여, 5\,L 참조수가 tank 위상 내내
동시 폭기되고 자기 차례에 빠르게 수렴하게 한다:
\begin{enumerate}[nosep,leftmargin=1.6em]
\item \emph{준비:} \texttt{m3b}(챔버에서 KCl 배출)
$\rightarrow$ \texttt{m1f}(본수조 $\to$ 홀딩)
$\rightarrow$ \texttt{m2f}(홀딩 $\to$ 챔버).
\item \emph{위상 A:} \texttt{ron}; 평탄까지 30\,s마다 \texttt{tank} 측정(챔버 tank +
5\,L 동시 폭기).
\item \emph{전이:} \texttt{m2b}로 tank 분취액을 홀딩에 파킹; \texttt{m4f}로 참조수를
챔버에 도입.
\item \emph{위상 B:} \texttt{ron}; 평탄까지 \texttt{ref} 측정(이미 평형 근처라 빠름).
\item \texttt{calkh}; 이후 정리: \texttt{m4b}(참조수 5\,L 회수), \texttt{m1b}(파킹된
tank 본수조 복귀), \texttt{m3f}(KCl 소크 복원), \texttt{airoff}(프로브 KCl에 안착).
\end{enumerate}

\subsection{이중 게이트 평탄 판정}
읽기값은 부동소수 경계 지터를 피하기 위해 정수 milli-pH(mpH)로 보관한다. 평탄은
슬라이딩 윈도우에서 \emph{두} 게이트가 모두 성립할 때만 선언된다(표~\ref{tab:plateau}):
짧은 \emph{span} 게이트(최근 4개, $\mathrm{max}-\mathrm{min}\le 2$\,mpH)는 잡음을
배제하고, 더 긴 \emph{net} 게이트(최근 8개, $|w_{-1}-w_{0}|\le 1$\,mpH)는 span
게이트만으로는 수렴으로 오인하는 느린 단조 지수 꼬리를 배제한다. 더 긴 net 룩백
(8 대 4; ``B1'')은 tank가 일찍 등반을 멈추는 조기 latch에 대한 해결책이었다. 위상별
벽시계ㆍ읽기수 상한($\mathtt{PHASE\_MAX}=7200\,\mathrm{s}$, $\mathtt{MEAS\_MAX}=240$)과
연속 실패 상한($\mathtt{FAIL\_MAX}=5$)이 루프를 묶어 무한 대기를 막는다.

\begin{table}[t]
\centering
\caption{평탄 추종 파라미터(\texttt{measure\_until\_flat}).}
\label{tab:plateau}
\footnotesize
\begin{tabular}{@{}llp{0.5\linewidth}@{}}
\toprule
상수 & 값 & 의미 \\
\midrule
\texttt{FLAT\_SPAN\_N}  & 4 & span 윈도우(지터 게이트) \\
\texttt{FLAT\_SPAN\_MPH}& 2 & 최근 4개: $\max-\min\le 2$\,mpH \\
\texttt{FLAT\_NET\_N}   & 8 & net 룩백(드리프트 게이트; B1) \\
\texttt{FLAT\_NET\_MPH} & 1 & 최근 8개: $|w_{-1}-w_0|\le 1$\,mpH \\
\texttt{MEAS\_INTERVAL} & 30\,s & 읽기 간격(폭기 지속) \\
\texttt{PHASE\_MAX\_SECS}& 7200\,s & 위상별 상한(2\,h; 총 4\,h) \\
\texttt{MEAS\_MAX}      & 240 & 위상별 읽기 상한 \\
\texttt{FAIL\_MAX}      & 5 & 허용 연속 파싱 실패 \\
\bottomrule
\end{tabular}
\end{table}

그림~\ref{fig:plateau}는 실측 접근 곡선 한 쌍이다. 갓 챔버에 들어온 tank는 느린
일차 꼬리로 천장을 향해 등반하고, net 게이트가 진짜 평탄이 될 때까지(8점 net
드리프트가 1\,mpH로 떨어지는 $\sim$624\,s에 latch) 붙잡는다. 반면 tank 위상 내내
5\,L에서 동시 폭기된 참조수는 이미 평형 상태라 $\sim$2.5\,min에 latch된다---tank를
먼저 측정하는 것이 빠른 이유이자, span 게이트만으로는($\sim$300\,s 부근에서 충족)
tank를 조기 latch했을 이유를 직접 보여 준다.

\begin{figure}[t]
\centering
\begin{tikzpicture}
\begin{axis}[
  width=0.86\linewidth, height=6.2cm,
  xlabel={경과 시간 (s)}, ylabel={pH (평형 접근)},
  xmin=-20,xmax=660, ymin=7.825, ymax=7.915,
  grid=both, major grid style={gray!18}, minor grid style={gray!8},
  legend style={at={(0.63,0.55)},anchor=center, fill=white, fill opacity=0.85,
    text opacity=1, draw=gray!40, font=\footnotesize}, legend cell align=left,
  tick label style={font=\footnotesize}, label style={font=\small},
]
\addplot[blue!70!black, mark=*, mark size=1.1pt, thick] coordinates {
(0,7.829)(38,7.832)(77,7.834)(116,7.836)(156,7.838)(195,7.840)(234,7.841)
(273,7.843)(312,7.844)(351,7.845)(390,7.846)(429,7.847)(468,7.848)(507,7.849)
(546,7.849)(585,7.850)(624,7.850)};
\addlegendentry{tank (챔버 시료)}
\addplot[red!75!black, mark=square*, mark size=1.1pt, thick, dashed] coordinates {
(0,7.906)(38,7.909)(77,7.909)(117,7.909)(156,7.909)};
\addlegendentry{ref (동시폭기 5\,L)}
\node[anchor=west,font=\scriptsize,blue!50!black] at (axis cs:430,7.843) {net 게이트 latch @624\,s};
\draw[blue!50!black,->,>=stealth] (axis cs:560,7.844) -- (axis cs:620,7.849);
\end{axis}
\end{tikzpicture}
\caption{실측 평탄 추종 곡선(밀폐, net 게이트 실행). tank 시료(실선)는 느린 지수
꼬리를 등반하며 이중 게이트 판정(최근 4개 span $\le 2$\,mpH \emph{및} 최근 8개 net
$\le 1$\,mpH)이 충족될 때만 latch된다; 동시폭기된 참조수(점선)는 이미 평형 근처라
빠르게 latch된다. 두 계열의 pH가 다른 것은 기저 시료가 다르기 때문이며(이 실행은 널
캠페인 이전), 각자의 \emph{평탄} 값만이 식~\eqref{eq:dkh}에 들어간다.}
\label{fig:plateau}
\end{figure}

\subsection{로깅과 오류 처리}
완료된 측정마다 \texttt{HH ref\_pH tank\_pH 8.448 tank\_kh temp} 한 줄이
\texttt{dkh.dat}에 추가된다. 두 가지 구분되는 표식이 비정상 결과를 인코딩한다:
하드 실패(타임아웃, 시리얼 손실, KCl 소크 미확인) 시의 \emph{전부 0} 줄---이는
추가로 \emph{래치}되어, 사람이 지울 때까지 이후 실행이 측정을 거부하고 표식만
재기록하여 무인 펌프 반복 구동으로부터 프로브를 보호한다; 그리고 위상이 평탄 미도달로
상한에 걸렸을 때의 \emph{음수} \texttt{tank\_kh}(크기는 보존, 부호만 반전)---서버가
이를 ``평탄 미도달''로 읽되 래치는 트리거하지 않는다. \texttt{finally} 블록은 임의의
비정상 종료 시 비상 정리(에어 오프, 챔버 배출, KCl 소크 복원)를 실행한다.

\subsection{널 테스트 프로토콜}
진단의 중추는 \emph{널 테스트}다: 참조수와 tank 용기를 \emph{같은} 물로 채워 진짜
알칼리니티 차를 0으로 만들면 정답은 정확히 $\dpH=0$, $\mathrm{dKH}=8.448$이다. 두
변형을 쓴다. \emph{부분 널}에서는 소부피 tank 분취액만 신선 해수로 교체하고 노화된
5\,L 참조수는 유지하여, \dpH{} 변화가 tank 물의 기여를 분리하게 한다. \emph{완전
널}에서는 tank 챔버에도 참조수를 넣어 물의 차이를 \emph{구성상} 0으로 만들어 남는
\dpH{}가 순수 계측 기인임이 되게 한다.

\subsection{개정 이력과 한 번에 하나씩 변경 원칙}
제작물은 5개 개정을 거쳤다(표~\ref{tab:versions}). 변경은 한 번에 하나의 인과 축으로만
이루어져 각 캠페인이 단일 메커니즘을 시험하게 했다; 두 상수를 함께 바꾼 경우
(예: 안정화 윈도우와 수렴 밴드)에는 둘이 두 개의 독립 변수가 아니라 \emph{같은}
인과 축(응답 완전성)임을 근거로 정당화하였다.

\begin{table}[t]
\centering
\caption{방법/펌웨어 개정. ``ron 유지''=측정 간 지속 폭기.}
\label{tab:versions}
\footnotesize
\begin{tabular}{@{}llp{0.56\linewidth}@{}}
\toprule
개정 & 기간 & 이전 대비 변경 \\
\midrule
V0 & 6/14--6/15 & 첫 대칭 제작; 참조수를 5\,L에서 컵으로 분취; 고정 1200\,s 폭기. \\
V1 & 6/15--6/16 & 측정 간 지속 폭기 추가(\texttt{ron} 유지). \\
V2 & 6/16--6/17 & 참조수를 소량 직접 폭기 분취로(홀딩 $\to$ 챔버); 측정 시 폭기 OFF. \\
V3 & 6/17 & 측정 \emph{중} 폭기(tank/ref를 같은 \pco{}에 고정); 경로 대칭화; KCl 헹굼 기포 청소(이후 폐기). \\
V4 & 6/17$\to$ & \textbf{평탄 추종}으로 진짜 평형까지(고정 시간 없음); 2-위상 tank-먼저-ref; 무한대기 방지 상한; 비상 정리. 이후: \emph{밀폐 공통 헤드스페이스}, \emph{5\,L 능동 앵커}, 이중 게이트(B1) 판정. \\
\bottomrule
\end{tabular}
\end{table}

% =====================================================================
\section{결과}\label{sec:results}
% =====================================================================
표~\ref{tab:data}가 전체 캠페인이다. 전 구간에서 $\dpH<0$은 tank가 참조수보다 더
산성으로 읽혔음을(겉보기 알칼리니티가 더 낮음) 뜻하며, 식~\eqref{eq:dkh}에 의해 원시
오차는 $\mathrm{dKH}-8.448 = 8.448\,(10^{\dpH}-1)$이다.

{\scriptsize\setlength{\tabcolsep}{4pt}
\begin{longtable}{@{}r l l r r r r r p{0.30\linewidth}@{}}
\caption{전체 측정 캠페인(29회). ``abs pH''는 공통모드 수준(하루 안에서만 실내 \co{}
프록시). 구성 약어: S=밀폐 헤드스페이스, A=5\,L 능동 앵커, N=net 게이트, B1=net 룩백 8.
\textbf{널 단계:} 시작 시 신선 해수를 한 배치 분취해 둔다; PN=부분 널($100\,\mathrm{mL}$
tank 챔버만 그 신선 분취 시료로 다시 채우고 노화된 $5\,\mathrm{L}$ 참조수는 유지 → tank
물의 기여 분리); CN=완전 널(이어서 참조수를 tank 챔버에도 옮겨 양 용기가 \emph{동일한}
물 → 물 차이 구성상 0).}
\label{tab:data}\\
\toprule
\# & 일시 & 구성 & pH$_\mathrm{ref}$ & pH$_\mathrm{tank}$ & \dpH{} & dKH & $T$ & 조건 \\
\midrule
\endfirsthead
\multicolumn{9}{l}{\scriptsize\emph{표~\ref{tab:data} (계속)}}\\
\toprule
\# & 일시 & 구성 & pH$_\mathrm{ref}$ & pH$_\mathrm{tank}$ & \dpH{} & dKH & $T$ & 조건 \\
\midrule
\endhead
\midrule
\multicolumn{9}{r}{\scriptsize\emph{다음 쪽에 계속}}\\
\endfoot
\bottomrule
\endlastfoot
1 & 06-14 20:12 & V0 & 7.959 & 8.021 & $+0.062$ & 9.741 & 26.7 & 첫 수동 대칭 측정 \\
2 & 06-14 21:00 & V0 & 7.997 & 8.072 & $+0.075$ & 10.054& 26.2 & --- \\
3 & 06-15 05:00 & V0 & 8.057 & 8.077 & $+0.020$ & 8.843 & 26.4 & 새벽인데 $+$; 창문 열림(환기), 저 \co{} \\
4 & 06-15 13:00 & V1 & 8.050 & 8.042 & $-0.008$ & 8.293 & 27.0 & 첫 지속 폭기 \\
5 & 06-15 21:00 & V1 & 8.109 & 8.091 & $-0.018$ & 8.109 & 27.3 & 깨끗한 저 \co{} 점 \\
6 & 06-16 05:00 & V1 & 7.810 & 7.751 & $-0.059$ & 7.362 & 27.5 & 새벽 닫힘방 \co{} 급등(양 abs pH $\sim$0.3 하락) \\
7 & 06-16 13:00 & V1b& 8.091 & 8.063 & $-0.028$ & 7.922 & 27.7 & 윈도우/밴드 강화 \\
-- & 06-16 21:00 & --- & \multicolumn{5}{c}{\emph{실패(0.000; 정전, 포트 세마포)}} & \\
8 & 06-16 21:21 & V2 & 8.042 & 8.021 & $-0.021$ & 8.039 & 28.6 & 첫 V2(수동) \\
9 & 06-17 05:00 & V2 & 7.997 & 7.939 & $-0.058$ & 7.386 & 27.8 & V2 새벽 검증 \emph{실패}; V1으로 회귀 \\
10& 06-17 13:00 & V3 & 8.109 & 8.091 & $-0.018$ & 8.095 & 28.2 & 첫 V3; ref 미평형(5 읽기) \\
11& 06-17 21:00 & V4 & 8.113 & 8.071 & $-0.042$ & 7.669 & --- & 첫 V4 + 절사평균 펌웨어; 깨끗한 평탄 \\
12& 06-18 05:00 & V4 & 7.738 & 7.695 & $-0.043$ & 7.657 & --- & \textbf{결정적:} 새벽 $-0.043\approx$ 저녁 $-0.042$, abs pH 0.4 스윙에도 $\Rightarrow$ \co{}-수준 무관 \\
13& 06-18 13:00 & V4 & 8.081 & 8.055 & $-0.026$ & 7.950 & 28.3 & \#11과 같은 저 \co{}인데 $-0.026$; dKH $+0.29$ 점프 $\Rightarrow$ baseline 비상수(증발, 개방 200\,mL 5일째) \\
14& 06-19 05:00 & V4+S & 7.976 & 7.968 & $-0.008$ & 8.283 & 28.1 & \textbf{밀폐 검증\,\textcircled{1}}: 새벽 dip 붕괴 $-0.05\!\to\!-0.008$; 평탄 33$\to$4 읽기 \\
15& 06-19 13:00 & V4+S & 7.957 & 7.947 & $-0.010$ & 8.255 & 28.6 & \textbf{밀폐 검증\,\textcircled{2}}: 새벽$\to$낮 점프 $+0.29\!\to\!-0.03$($\sim$10$\times$); 증발 제거 \\
16& 06-19 21:00 & V4+S & 7.969 & 7.968 & $-0.001$ & 8.421 & 26.8 & 거의 완벽한 널(원시오차 $-0.027$, 최소); 단 \#14/15 대비 $+0.166$ $\Rightarrow$ offset 완벽 상수 아님; 냉방 과도($-1.8$\,\textcelsius) \\
17& 06-20 05:00 & V4+S & 7.922 & 7.871 & $-0.051$ & 7.521 & 27.7 & 새벽 dip 재출현; tank 11$\times$ 단조=조기 latch(진짜 평탄 아님); 냉방이 kinetics 둔화 \\
18& 06-20 15:00 & V4+SAN& 7.975 & 7.929 & $-0.046$ & 7.602 & 27.2 & 첫 5\,L 앵커; tank \emph{진짜} 평탄(34$\times$, 21.6\,min)인데도 $-0.046$ $\Rightarrow$ \textbf{실제 gap}, latch 아티팩트 아님 \\
19& 06-20 18:11 & V4+S,2$\times$& 7.979 & 7.924 & $-0.055$ & 7.443 & 26.4 & 통제실험; 2$\times$ 폭기가 등반 반감(1296$\to$711\,s)하나 \dpH{} 불변 $\Rightarrow$ kinetics \emph{아님} \\
20& 06-20 18:57 & V4+S,2$\times$& 7.975 & 7.931 & $-0.044$ & 7.636 & 26.7 & ref-먼저 + 호스 스왑; gap 지속 $\Rightarrow$ 순서/경로/미평형 모두 기각 \\
21& 06-20 19:50 & V4+S & 7.975 & 7.926 & $-0.049$ & 7.539 & 26.9 & 재현 \\
22& 06-20 20:29 & V4+SA& 7.977 & 7.921 & $-0.056$ & 7.426 & 26.9 & tank-먼저+파킹; ref 역대 최속(4$\times$); tank 조기 latch(7.921) $\Rightarrow$ B1 동기 \\
23& 06-20 21:00 & V4+SA,B1& 7.979 & 7.926 & $-0.053$ & 7.474 & 27.0 & \textbf{B1 성공}: tank 25$\times$/945\,s 진짜 평탄; \#22가 $\sim$5\,mpH 미달이었음 확인; gap(7번째) \\
24& 06-21 05:00 & V4+SA,B1& 7.974 & 7.928 & $-0.046$ & 7.594 & 26.7 & 새벽이 저녁 군집과 일치 $\Rightarrow$ 시간대 의존성 배제(gap, 8번째) \\
25& 06-21 06:00 & V4+S,PN& 7.973 & 7.963 & $-0.010$ & $-8.255$ & 26.7 & \textbf{부분 널}: tank만 신선; tank pH $+37$\,mpH, \dpH{} $-0.05\!\to\!-0.010$; tank 미평탄(옛 상한) $\Rightarrow$ 하한값 \\
26& 06-21 13:00 & V4+S,PN& 7.964 & 7.956 & $-0.008$ & 8.298 & 27.2 & \textbf{부분 널, 완전평형}: dKH 8.298, 원시오차 $-0.15$; 좋은 군집 복귀; 4\,h 상한 검증 \\
27& 06-21 19:00 & V4+S,CN& 7.960 & 7.957 & $-0.003$ & 8.387 & 27.4 & \textbf{완전 널}: 참조수를 tank 챔버에도 $\Rightarrow$ 양쪽 동일; 원시오차 $-0.061$=순수 계측 오프셋; 100\,mL 대 5\,L 상쇄 \\
28& 06-21 21:00 & V4+S,CN& 7.959 & 7.958 & $-0.001$ & 8.434 & 27.4 & \textbf{완전 널}(2번째): $\dpH-0.001$(분해능 한계), 원시오차 $-0.014$(최소); $\le0.06\dkh$ 오프셋 재확인 \\
-- & 06-22 05:00 & V4+S,CN & \multicolumn{5}{c}{\emph{실패(상한)}} & 55분 스케줄러 상한이 ref 도중 강제종료; tank는 평탄 완료(7.955); 해법 5\,h 상한 \\
29& 06-22 09:00 & V4+S,CN& 7.959 & 7.948 & $-0.011$ & 8.226 & 27.4 & \textbf{노화 완전 널}($\sim$14\,h, $100\,\mathrm{mL}$ 노출 분취액): 널이 원시오차 $-0.22$로 붕괴; tank 진짜 평탄 59$\times$(긴 폭기가 석출 유발) \\
\bottomrule
\end{longtable}
}

\subsection{새벽 \co{} 오차와 그 진단}
가장 먼저 분명해진 병리는 \dpH{}의 새벽 붕괴였다(\#6, $-0.059$; 오차 $-1.09\dkh$).
결정적 단서는 공통모드였다: 밤사이 \emph{양} 절대 pH가 $\sim0.3$ 떨어져(ref
$8.109\!\to\!7.810$, tank $8.091\!\to\!7.751$) 낮의 상승 추세를 뒤집은 것---폭기수를
산성화하는 닫힌방 \co{} 축적의 지문이다. 자연 실험이 이를 확인했다: 환기한 단 한 밤
(\#3, 창문 열림)에 절대 pH가 높게 유지되고 \dpH{}는 정상인 $+0.020$이었던 반면, 닫힌
밤들은 $-0.058/-0.059$였다. 중요하게도, 실내 \co{}는 차동이 \emph{상쇄해야 할} 바로
그 공통모드 변수다; 그것이 \dpH{}로 누출되었다는 것은 환기로 없앨 환경 문제가 아니라
ref/tank 대칭이 깨졌다는 신호였다.

\subsection{평탄 추종과 위상 간 드리프트 충돌}
진짜 평탄까지 측정(V4)하니 \co{}-\emph{수준} 의존성이 사라졌다: \#11/\#12가 절대
pH 0.4 스윙을 가로질러 $\dpH=-0.042$, $-0.043$을 주었다. 그러나 첫 V4 실행은 경쟁
실패를 드러냈다: 31분짜리 tank 측정에서 pH가 정점까지 올랐다 \emph{내려가}(측정 중
실내 \co{} 드리프트), 단일 프로브가 각 용기를 우연히 샘플링한 시점에 따라 정해지는
값에서 latch되었다. 이 \emph{위상 간 드리프트}는 대략
\begin{equation}
\epsilon_\mathrm{drift} \;\approx\; \Delta t_{\mathrm{tank{-}ref}}\times \frac{dC}{dt},
\label{eq:drift}
\end{equation}
이며, 단일 프로브에서 근본 충돌을 일으킨다: 완전 평형은 긴 측정을 요하지만, 긴
측정은 tank--ref 샘플링 간격 $\Delta t$를 넓혀 드리프트를 키운다. $dT/dt$를 쓴 같은
형태가 \#16을 설명한다---냉방 과냉($-1.8$\,\textcelsius{} 측정 중)이 $+0.166\dkh$ 변동과
일치했다.

\subsection{밀폐가 두 교란을 붕괴시킨다}
단일 \emph{밀폐 공통 헤드스페이스}가 이 충돌을 원천에서 해결한다. ref와 tank가 하나의
공기 상자를 공유하게 강제함으로써, 공통 \pco{}를 우연이 아닌 본질로 만들고, 실내
\co{} 변동을 차단하며, (보너스로) 개방 tank 분취액을 농축하던 증발을 억제한다. 효과는
즉각적이었다(\#14/\#15): 새벽 dip이 $-0.05$에서 $-0.008$로 붕괴; 새벽$\to$낮 dKH
점프가 $\sim$10$\times$ 축소(개방 $+0.29$ $\to$ 밀폐 $-0.03$); 평탄 시간이 $\sim$33
읽기($\sim$50\,min)에서 $\sim$4($\sim$5\,min)로 격감---이것이 다시
식~\eqref{eq:drift}의 $\Delta t$를 붕괴시켜 위상 간 드리프트를 구조적으로 제거한다.
무보정 오차는 $-0.78$에서 $\sim-0.18\dkh$로 떨어졌다.

\subsection{능동 앵커와 실제 gap의 지속}
밀폐만으로는 느린 헤드스페이스 드리프트가 남았다. 수동 공유 공기 부피는 표면 확산
(수십 분)으로만 재평형하기 때문이다. 참조수 \emph{안에} 에어스톤을 두면(``앵커'')
공유 \pco{}를 수 분 내 능동 고정한다; 부피 감사가 지배 위계를 확인한다(시료
$\sim$3\,\si{\micro mol} $\ll$ 헤드스페이스 $\sim$65\,\si{\micro mol} $\ll$ 5\,L
$\sim$500\,\si{\micro mol}). 그런데 앵커와 net 게이트를 갖추고도, tank가 검증된 진짜
평탄에 도달했을 때조차 $\sim$$-0.046$--$-0.053$ gap이 \emph{지속}되었다(\#18: 34
읽기, 21.6\,min). 이어진 통제실험 쌍(\#19/\#20)이 네 가지 장치 측 설명을 차례로
기각했다: 폭기를 2배로 하면 등반 시간은 반감했으나 \dpH{}는 불변(kinetics \emph{아님});
측정 순서를 뒤집고 호스를 스왑해도 부호가 뒤집히지 않음(순서/웜업/유로 \emph{아님});
head-start한 tank가 \emph{같은} 평형에 안착(미평형 \emph{아님}). gap은 \emph{물}의
성질이었다.

\subsection{널 실험이 오차를 국소화한다}
장치 가설이 소진되었으므로 물을 직접 시험했다. \emph{부분 널}(\#25, tank 챔버만 신선
해수)이 \dpH{}를 $-0.05$에서 $-0.010$로 붕괴시켰다(tank pH $+37$\,mpH); 완전평형
(\#26)에서는 $\mathrm{dKH}=8.298$, 원시오차 $-0.15$---좋은 군집으로 복귀. tank에
신선 물을 넣자 거의 전 오차가 회복되어 노화된 tank 분취액이 지목되었다. \emph{완전
널}(\#27), \emph{양} 용기에 참조수를 넣어 물 차이를 구성상 0으로 만든 경우,
$\dpH=-0.003$, $\mathrm{dKH}=8.387$, 원시오차 $\mathbf{-0.061}$을 주었고,
$100\,\mathrm{mL}$:$5\,\mathrm{L}$ 부피 불일치에도 양쪽이 검증된 진짜 평탄이었다---부피
무관성의 직접 확인(\ref{sec:principle}절).

\subsection{오차 분해와 정확도}
두 널이 헤드라인 오차를 분해한다(표~\ref{tab:budget}, 그림~\ref{fig:budget}):
\begin{equation}
\underbrace{-0.92\dkh}_{\text{원시}} \;=\;
\underbrace{-0.86\dkh}_{\substack{\text{노화 tank 알칼리니티 손실}\\(\ce{CaCO3}\text{ 석출/생물막})}}
\;+\;
\underbrace{-0.06\dkh}_{\substack{\text{계측 오프셋}\\(\dpH=-3\,\mathrm{mpH})}}.
\end{equation}
7점 군집(\#18--24)에서 \emph{정밀도}는 $\dpH = -0.0499\pm0.0048$($\approx$5\,mpH),
$\mathrm{dKH}=7.531\pm0.084$---목표 $\pm0.25\dkh$의 약 1/3 산포다. 대부분의 오차가
시료 아티팩트로 규명되고 계측 오프셋이 $-0.06\dkh$(목표의 1/4)에 머물므로, 본
모니터는 \emph{무보정}으로 $\pm0.25\dkh$를 충족한다; 영점 보정은 불필요로 판정되었고,
정상(비열화) 산호수라면 무보정 오차는 $\pm0.1\dkh$ 안에 들 것이다.

\begin{table}[t]
\centering
\caption{전체 진단 후 오차 예산.}
\label{tab:budget}
\footnotesize
\begin{tabular}{@{}lcl@{}}
\toprule
성분 & 크기 & 성격 / 해법 \\
\midrule
노화 tank 알칼리니티 손실 & $-0.86\dkh$ & 시료 취급; 신선 분취액 사용 \\
계측 오프셋 & $-0.06\dkh$ & $\dpH=-3$\,mpH; 보정 가치 미만 \\
정밀도(1$\sigma$) & $\pm0.08\dkh$ & 전극 잡음 + 잔여 드리프트 \\
실내 \co{} 의존성 & $\approx 0$ & 밀폐로 제거 \\
증발 교란 & $\approx 0$ & 밀폐로 제거 \\
\bottomrule
\end{tabular}
\end{table}

\begin{figure}[t]
\centering
\begin{tikzpicture}
\begin{axis}[
  ybar, bar width=30pt,
  width=0.82\linewidth, height=5.8cm,
  symbolic x coords={Raw,Degrade,Offset},
  xtick=data,
  xticklabels={원시 오차,노화 tank 열화,계측 오프셋},
  nodes near coords, nodes near coords align={below},
  every node near coord/.append style={font=\footnotesize,
    /pgf/number format/.cd, fixed, fixed zerofill, precision=2},
  ymin=-1.05, ymax=0.12,
  ylabel={오차 기여 (dKH)},
  tick label style={font=\footnotesize}, label style={font=\small},
  ymajorgrids, major grid style={gray!18},
  enlarge x limits=0.28,
]
\addplot[draw=black, fill=blue!22] coordinates {(Raw,-0.92)(Degrade,-0.86)(Offset,-0.06)};
\draw[gray] (axis cs:Raw,0) -- (axis cs:Offset,0);
\end{axis}
\end{tikzpicture}
\caption{$-0.92\dkh$ 원시 오차의 분해. 완전 널이 계측 오프셋을 $-0.06\dkh$(목표
$\pm0.25\dkh$의 1/4)로 격리하고; 부분 널이 남은 $-0.86\dkh$를 노화된 tank 분취액의
알칼리니티 손실(\ce{CaCO3} 석출/생물막)로 귀속시킨다---측정 결함이 아니라 시료
아티팩트다.}
\label{fig:budget}
\end{figure}

\subsection{널의 시간 의존적 붕괴}
이 분해는 반증 가능한 예측을 낳는다: 대부분의 오차가 전자장치가 아니라 시료 화학이라면,
완벽해야 할 널이 시료가 노화될수록 \emph{쇠퇴}해야 한다. 우리는 같은 $100\,\mathrm{mL}$
분취액으로 다음 날 완전 널을 재측정하여(장치는 손대지 않음) 이를 시험했다
(그림~\ref{fig:degrade}, \#29). 원시 오차는 0\,h와 2\,h에 $-0.061$, $-0.014\dkh$였다가
$\sim$14\,h에 $-0.222\dkh$로 떨어졌다($\dpH=-0.011$; tank pH $7.957\!\to\!7.948$,
더 큰 참조수는 $7.959$ 유지). 이 하락은 $\pm0.08\dkh$ 정밀도를 훨씬 넘고 검증된 진짜
평탄(tank 59회)에서 도달했으므로, latch 아티팩트가 아니라 진짜 알칼리니티 손실이다.
소부피ㆍ노출 분취액과 비정상적으로 긴 폭기---탄산칼슘 과포화를 높여 \emph{측정 중}
\ce{CaCO3} 석출을 일으킴---는 $-0.86\dkh$ 메커니즘의 worst-case이며, 그에 따라 널이
무너진 반면 계측 오프셋 바닥은 신선 값을 유지했다. 이는 오차 분해의 시간 영역 확증이자
실무 경고다: 소부피 시료에서 비정상적으로 긴 평탄은 그 자체가 열화 신호다.

\begin{figure}[t]
\centering
\begin{tikzpicture}
\begin{axis}[
  width=0.82\linewidth, height=5.6cm,
  xlabel={첫 완전 널 이후 경과 시간 (h)},
  ylabel={원시 널 오차 (dKH)},
  xmin=-1, xmax=15.5, ymin=-0.27, ymax=0.06,
  grid=both, major grid style={gray!18},
  tick label style={font=\footnotesize}, label style={font=\small},
]
\draw[gray,thick] (axis cs:-1,0) -- (axis cs:15.5,0);
\draw[green!45!black,dashed] (axis cs:-1,-0.25) -- (axis cs:15.5,-0.25);
\node[anchor=south west,font=\scriptsize,green!40!black] at (axis cs:-0.8,-0.245) {$\pm0.25\dkh$ 목표 한계};
\addplot[blue!70!black, mark=*, thick,
  error bars/.cd, y dir=both, y explicit]
  coordinates {
  (0,-0.061) +- (0,0.08)
  (2,-0.014) +- (0,0.08)
  (14,-0.222) +- (0,0.08)
  };
\node[anchor=north east,font=\scriptsize,blue!45!black,align=right] at (axis cs:13.6,-0.10)
  {$100\,$mL 분취액,\\ $\sim$14\,h, 긴 폭기};
\end{axis}
\end{tikzpicture}
\caption{완전 널의 시간 의존적 붕괴. \emph{양} 용기에 참조수를 넣어(물 차이 구성상 0)
원시 오차는 처음 몇 시간 동안 $\pm0.08\dkh$ 정밀도(오차 막대)와 $\pm0.25\dkh$ 목표
안에 머물다가, 소부피($100\,\mathrm{mL}$)ㆍ노출 tank 분취액이 $\sim$14\,h 노화된 뒤
$-0.22\dkh$로 떨어진다---긴 평탄 추종 실행(59회)이 폭기 중 \ce{CaCO3} 석출을 더 유발한다.
전자장치는 그대로이고 시료 알칼리니티만 드리프트하므로, 지배 오차항이 시료 화학임을
직접 확증한다(그림~\ref{fig:budget} 참조).}
\label{fig:degrade}
\end{figure}

% =====================================================================
\section{고찰}
% =====================================================================
\paragraph{진단이 얻어낸 것.} 안정된 $\sim-0.8\dkh$ 오프셋을 보면 영점 보정하려는
충동이 든다. 이를 자제한 것---보정 상수는 독립 재현 후에만, 그리고 작은 비독립 표본에는
결코 적용하지 않는다는 방법론 원칙에 따라---이 결정적이었다. 그 오프셋은 알고 보니
상수인 장치 성질이 \emph{아니었고}(밀폐 군집 전반에서 $\pm0.085\dkh$ 변동했고 결국
대부분이 열화된 시료수로 판명), 그것을 보정으로 지웠다면 시료 아티팩트를 장치에
각인시켜 진짜 $-0.06\dkh$ 바닥을 가렸을 것이다.

\paragraph{외부 검증.} 우리가 독립적으로 도출한 오차 구조는 상용 AquaWiz의 보고된
거동을 그대로 반영하여 외부 지지를 받는다. 커뮤니티 보고는 실내 \co{} 변동을 지배
오차로, 밀폐 용기를 현장 해법으로 지목한다; 제조사 자체 안내---급격한 \co{} 변화는
``약 한 시간만 측정에 영향''을 주며 ``공기 흡입관을 스키머의 실외 공기 라인에
연결''하라---는 우리의 $\sim$1시간 위상 간 드리프트 에피소드와 가스 기준 고정 해법
양쪽에 부합한다(우리는 실외 라인 대신 밀폐로 구현). 두 번째로 많이 보고된 상용 오차인
염 침착에 의한 기포 품질 저하($>2\dkh$ 오차, 6개월마다 청소)는 우리의 폭기 충분성
분석과 일치한다. 보고된 현장 정확도(Hanna 대비 $\le0.2\dkh$; 무보정 7주 $\le0.75\dkh$;
``$0.5\dkh$ 이하=훌륭, 추세용'')는 우리의 $\pm0.25\dkh$ 결과와 추세 지향 의도를
괄호친다. 특히 상용 장치는 보정과 청소로 관리할 뿐 진짜 평형 평탄을 추구하지 않는다;
우리의 평탄 추종 단계가 한 발 더 나간다.

\paragraph{과학적 근거.} 방법의 기반인 동역학 분리는 잘 확립되어 있다: 수용성 \co{}
수화/탈수 완화는 $\sim$16\,s에 완료되므로\cite{johnson1982} 화학은 결코 율속이 아니다;
느린 단계는 물리 가스 교환으로, 직접 버블링에서 $\sim$15\,min이며\cite{takahashi}
우리가 관찰한 평탄 시간과 정합한다. ``\pco{}-평형화와 pH로부터 \at{}''라는 차동 원리는
$0.5\,\mathrm{mL}$ 소부피까지 $\pm1$--$2\,\si{\micro mol\,kg^{-1}}$ 정밀도로 입증된
최근 분석화학 근거가 있어\cite{fleger2025}, 소부피 소형화가 옳은 방향임을 지지한다.

\paragraph{한계.} (i) 단일 프로브에서는 식~\eqref{eq:drift}의 위상 간 드리프트를 0으로
몰 수 없고 억제만 가능하다(밀폐와 $\Delta t$ 최소화로). (ii) 역사적 지배 오차가 시료
열화였으므로, 운영 규율---신선하고 충분한 크기의 tank 분취액을 쓰고 보관 대신
순환---은 부가물이 아니라 장치의 일부다. (iii) 느린 드리프트 하의 긴 평탄은 측정을 수
시간으로 늘릴 수 있다; 벽시계 상한이 이를 묶지만 외부 스케줄러 한도를 코드의 4시간
천장 위로 설정해야 한다. (iv) 절대 pH는 하루 안 \co{} 프록시로만 유의하며, \dpH{}만이
정량적이다.

\paragraph{권장 운영.} 공통 헤드스페이스를 밀폐하고; 참조수를 능동 폭기로 앵커하며;
이중 게이트 평탄 판정으로 tank-먼저-ref를 측정하고; 신선한 시료수를 쓰며; 에어 경로를
주기적으로 청소하고; 출력은 추세 계기로 다루되, 보정한다면 ref$\leftrightarrow$tank
괴리에 연동된 월 주기로 독립 키트 대비 영점만 잡는다.

% =====================================================================
\section{결론}
% =====================================================================
무시약, 단일프로브 차동 pH 폭기 모니터는, 세 설계 규칙---단일 밀폐 공통 헤드스페이스,
진짜 가스 평형까지의 평탄 추종, 능동 폭기 참조수 앵커---을 지키면 영점 보정 \emph{없이}
산호 탄산알칼리니티를 $\pm0.25\dkh$보다 잘 측정할 수 있다. 9일 캠페인의 가치는 어떤
한 수치보다 $-0.92\dkh$ 원시 오차를 \emph{분해}한 데 있었다: 널 실험은 그중 $-0.86\dkh$가
열화된 시료 분취액의 알칼리니티 손실이고 단 $-0.06\dkh$만이 장치임을 보였다. 차동 설계는
실내 \co{}와 프로브 드리프트의 공통모드 제거라는 약속을 이행했고; 남은 것은 전자장치의
오차가 아니라 비커 속 화학이었다. 우리는 이 설계 규칙과 널 분해 방법론이 다른 저가
알칼리니티 계측에 이전되리라 기대한다.

% ---------------------------------------------------------------
\section*{재현성}
펌웨어, 호스트 측정 스크립트(\texttt{measure\_kh\_once.py}), 원시 로그
(\texttt{dkh.dat}), 전체 측정 대장은 프로젝트 저장소~\cite{repo}에서 이용할 수 있다.

\section*{감사의 글}
REEF FAMILY 망고미역님과 DIY 모임 참여자 모두의 열정에 힘입어 이 프로젝트를 마무리
할 수 있게 되었습니다.

% ---------------------------------------------------------------
\begin{thebibliography}{9}
\bibitem{repo} nowforyou (REEF FAMILY), ``ReefKeeper: 개방형 차동 pH 폭기
알칼리니티 모니터(펌웨어, 호스트 스크립트, 측정 대장),'' 프로젝트 저장소,
\url{https://github.com/taeseokyi/reefkeeper}.
\bibitem{johnson1982} K.~S.~Johnson, ``Carbon dioxide hydration and dehydration
kinetics in seawater,'' \emph{Limnol.\ Oceanogr.} \textbf{27}(5), 849--855
(1982). doi:10.4319/lo.1982.27.5.0849.
\bibitem{soli} A.~L.~Soli and R.~H.~Byrne, ``CO\textsubscript{2} system
hydration and dehydration kinetics and the equilibrium CO\textsubscript{2}/H\textsubscript{2}CO\textsubscript{3}
ratio in aqueous NaCl solution,'' \emph{Mar.\ Chem.} (2002).
\bibitem{zeebe2001} R.~E.~Zeebe and D.~Wolf-Gladrow, \emph{CO\textsubscript{2}
in Seawater: Equilibrium, Kinetics, Isotopes}, Elsevier Oceanography Series 65
(2001).
\bibitem{takahashi} T.~Takahashi, ``Methods for measurement of
$p\mathrm{CO_2}$ and total CO\textsubscript{2} in seawater'' (탄산계 방법론;
직접 버블링 평형화 $\sim$15\,min).
\bibitem{fleger2025} K.~Fleger, X.~Liu, W.~Berelson, J.~Adkins \emph{et al.},
``Total alkalinity measurements in small samples based on
CO\textsubscript{2} equilibration and spectrophotometric pH,''
\emph{Anal.\ Chim.\ Acta} (2025). doi:10.1016/j.aca.2025.344432.
\bibitem{roberts2005} (가스 스트리핑 일차 평형화;
$\ln(C_0/C_t)\propto \phi/V$.) Roberts, \emph{Geophys.\ Res.\ Lett.} (2005).
\bibitem{egleston2010} C.~S.~Egleston, C.~L.~Sabine, F.~M.~M.~Morel,
``Revelle (완충) 인자, 해수 탄산계,''
\emph{Global Biogeochem.\ Cycles} (2010).
\bibitem{community} 상용 AquaWiz 알칼리니티 모니터에 대한 산호 취미가 커뮤니티 현장
보고(실내 \co{} 오차, 밀폐 용기 해법, 기포 저하, 제조사 안내): Reef2Reef 스레드
\#1159272, \#1073924, \#1115439; UltimateReef \#939234; Bay Area Reefers \#40433.
\end{thebibliography}

\end{document}
